\chapter{Data Frame Fundamentals}
\label{chap:data-frames}

\chapterguide
  {You should understand vectors, lists, indexing, and basic control flow.}
  {create, inspect, index, extend, and reduce data frames.}

\section{Creating a data frame}

Create an employee table from equal-length vectors:

\begin{lstlisting}[style=dsr]
emp.data <- data.frame(
  emp_id = 1:5,
  emp_name = c("Rick", "Dan", "Michelle", "Ryan", "Gary"),
  salary = c(623.3, 515.2, 611.0, 729.0, 843.25)
)
\end{lstlisting}

A data frame stores named, equal-length columns that may have different types.

\section{Creating a data frame from another structure}

\subsection{From a matrix}

\begin{lstlisting}[style=dsr]
my_matrix <- matrix(c(1, 2, 3, 4, 5, 6), nrow = 2)
print(my_matrix)

df_from_matrix <- data.frame(my_matrix)
print(df_from_matrix)

names(df_from_matrix) <- c("col_1", "col_2", "col_3")
print(df_from_matrix)
\end{lstlisting}

The generated column names are then replaced.

\subsection{From a list of vectors}

\begin{lstlisting}[style=dsr]
my_list <- list(
  rating = 1:4,
  animal = c("koala", "hedgehog", "sloth", "panda"),
  country = c("Australia", "Italy", "Peru", "China"),
  avg_sleep_hours = c(21, 18, 17, 10)
)

print(my_list)
super_sleepers <- data.frame(my_list)
print(super_sleepers)
\end{lstlisting}

The equal-length list components become columns.

\section{Displaying and inspecting a data frame}

Print the table directly:

\begin{lstlisting}[style=dsr]
print(emp.data)
\end{lstlisting}

RStudio's viewer displays the same object in a spreadsheet-like window.

\subsection{Structure}

\begin{lstlisting}[style=dsr]
str(emp.data)
\end{lstlisting}

\texttt{str()} reports columns and stored types.

\subsection{Summary and dimensions}

Use summary and dimension functions for different purposes:

\begin{lstlisting}[style=dsr]
print(summary(emp.data))

print(dim(emp.data))
print(ncol(emp.data))
print(nrow(emp.data))
\end{lstlisting}

\texttt{summary()} describes columns; \texttt{dim()}, \texttt{ncol()}, and \texttt{nrow()} report shape.

\begin{pitfall}
Lab 6a appears to swap the headings for summary and dimension output. The functions above are organized by their actual behavior.
\end{pitfall}

\section{Accessing columns and rows}

The salary column is extracted in three forms:

\begin{lstlisting}[style=dsr]
print(emp.data$salary)
print(emp.data[["salary"]])
print(emp.data[3])
\end{lstlisting}

\texttt{\$} and \texttt{[[} return a column vector; \texttt{[3]} returns a one-column data frame.

The first two rows with all columns are selected using:

\begin{lstlisting}[style=dsr]
result <- emp.data[1:2, ]
print(result)
\end{lstlisting}

Data frames use \texttt{[row, column]} indexing.

\section{Adding columns}

Add a column by assigning a new name:

\begin{lstlisting}[style=dsr]
emp.data$dept <- c("IT", "Operations", "IT", "HR", "Finance")
print(emp.data)
\end{lstlisting}

\texttt{cbind()} can attach another column:

\begin{lstlisting}[style=dsr]
emp.newdata <- cbind(
  emp.data,
  start_date = as.Date(c(
    "2012-01-01",
    "2013-09-23",
    "2014-11-15",
    "2014-05-11",
    "2015-03-27"
  ))
)
\end{lstlisting}

The resulting table includes character, numeric, and date columns.

\section{Adding rows}

Append compatible rows with \texttt{rbind()}:

\begin{lstlisting}[style=dsr]
emp.newdata <- data.frame(
  emp_id = c(6:8),
  emp_name = c("Rasmi", "Pranab", "Tusar"),
  salary = c(578.0, 722.5, 632.8),
  dept = c("IT", "Operations", "Finance")
)

emp.finaldata <- rbind(emp.data, emp.newdata)
print(emp.finaldata)
\end{lstlisting}

The column names and compatible types must match.

\section{Removing rows and columns}

Negative indices exclude rows or columns:

\begin{lstlisting}[style=dsr]
# Remove the second row and first column
result <- emp.data[-2, -1]
print(result)
\end{lstlisting}


\section{Extra practice table}

Extend this practice table with additional rows, columns, and types:

\begin{lstlisting}[style=dsr]
friends <- data.frame(
  friend_id = 1:5,
  friend_name = c("Sachin", "Sourav", "Dravid", "Sehwag", "Dhoni"),
  location = c("Kolkata", "Delhi", "Bangalore", "Hyderabad", "Chennai")
)
print(friends)
\end{lstlisting}

\begin{quickcheck}
\begin{enumerate}
  \item What does \texttt{str(emp.data)} help you inspect?
  \item Which functions in the handout report the number of rows and columns?
  \item Give two supplied ways to access the salary column by name.
  \item Which function binds new columns and which function binds new rows?
  \item How is negative indexing reused to remove data-frame rows or columns?
\end{enumerate}
\end{quickcheck}

\begin{chapterrecap}
\begin{itemize}
  \item Data frames combine equal-length columns of potentially different types.
  \item Inspect structure, summaries, and dimensions with dedicated functions.
  \item Index with \texttt{\$}, brackets, or row/column coordinates; extend with \texttt{cbind()} and \texttt{rbind()}.
\end{itemize}
\end{chapterrecap}
